\documentclass[11pt]{article}
\usepackage[margin=1in]{geometry}
\usepackage{amsmath,amssymb,amsthm}
\usepackage{bm}
\usepackage{enumitem}
\usepackage[hidelinks]{hyperref}
\usepackage{booktabs}

\newtheorem{definition}{Definition}
\newtheorem{remark}{Remark}
\newtheorem{property}{Property}

% ==== convenience macros ====
\newcommand{\R}{\mathbb{R}}
\newcommand{\E}{\mathbb{E}}
\newcommand{\Fspace}{\mathcal{F}}
\newcommand{\gates}{\mathcal{G}}
\newcommand{\proto}{\Theta}
\newcommand{\evalop}{\mathcal{E}}
\newcommand{\feasible}{\Fspace_{\proto}}
\newcommand{\sign}{\mathbf{z}}
\newcommand{\wt}{\mathbf{w}}
\newcommand{\zoo}{\Fspace_{\mathrm{zoo}}}
\newcommand{\pred}{\hat{\bm{y}}}
\newcommand{\combiner}{\mathcal{M}}

\title{\vspace{-2.5em}\textbf{Problem Formulation (Section 3)}\\[0.3em]
\large A Strict, Net-of-Cost, Capacity-Aware Objective for Agentic Alpha Mining}
\date{}

\begin{document}
\maketitle
\vspace{-3em}

%======================================================================
\section{Problem Formulation}
\label{sec:problem}

We consider the formulaic alpha-mining setting, in which an automated, multi-agent
generation process proposes candidate factors that are evaluated on historical market
data. Our contribution in this work is confined to the \emph{objective}: what a
candidate factor is scored on, and under what protocol that score is computed. We first
restate the frictionless predictive objective that prevails in the LLM-driven
factor-mining literature~\cite{han2026quantaalpha,shi2025alphajungle}
(\S\ref{sub:prelim}), make explicit the gap
between that objective and tradeable performance (\S\ref{sub:gap}), and then define a
deterministic execution model with transaction cost \emph{and} latency
(\S\ref{sub:exec}), a multi-objective factor-quality vector (\S\ref{sub:objectives}),
and a strict, non-modifiable backtesting protocol (\S\ref{sub:protocol}). The full
problem statement is given in \S\ref{sub:statement}.

%----------------------------------------------------------------------
\subsection{Preliminaries and the frictionless objective}
\label{sub:prelim}

Let $\mathcal{S}=\{s_1,\dots,s_N\}$ be a cross-section of $N$ tradeable assets and
$\mathcal{T}=\{t_1,\dots,t_T\}$ a sequence of trading periods. Market observations
form a feature tensor $\bm{X}\in\R^{N\times T\times D}$ over $D$ point-in-time
features (e.g., open, high, low, close, volume, VWAP), with $\bm{X}_t\in\R^{N\times D}$
the cross-sectional slice available at the close of period $t$. An \emph{alpha factor}
is a function $f\in\Fspace$ mapping the information available at $t$ to a
cross-sectional signal
\begin{equation}
  f(\bm{X}_{\le t}) \;=\; \sign_t \in \R^{N},
  \label{eq:factor}
\end{equation}
intended to be predictive of the forward cross-sectional return
$\bm{y}_{t+1}\in\R^{N}$. The base label is the next-period close-to-close return with a
one-period implementation gap,
$y_{i,t}=P^{\mathrm{close}}_{i,t+2}/P^{\mathrm{close}}_{i,t+1}-1$, a standard choice in
the daily factor-mining literature; we generalize this gap into an explicit latency
parameter in \S\ref{sub:exec}.

\paragraph{Factor combination.}
Individual factors are not traded in isolation. An \emph{effective-alpha repository}
$\zoo\subseteq\Fspace$ accumulates the factors that have passed evaluation, and a
supervised combination model $\combiner$ maps the vector of their signals to a single
composite return prediction,
\begin{equation}
  \pred_t \;=\; \combiner\big(\,[\,f(\bm{X}_{\le t})\,]_{f\in\zoo}\,;\,\theta\,\big)\in\R^{N},
  \label{eq:combiner}
\end{equation}
where $\combiner$ is a gradient-boosted decision-tree regressor (LightGBM). Its
parameters $\theta$ are fit by supervised
regression on the training split only, taking as inputs the signals of the factors in
the \emph{frozen} library $\zoo$ and as target the realized label; the library is held
fixed for each refit of $\combiner$, and no out-of-sample or test data enters the fit
(\S\ref{sub:protocol}). The portfolio in \S\ref{sub:exec} is constructed from the
composite prediction $\pred_t$, whereas the per-factor quality vector of
\S\ref{sub:objectives} is computed on the individual signal $\sign_t$; a new factor is
thus judged both on its stand-alone merit and, through Eq.~\eqref{eq:combiner}, on its
marginal contribution to the combined book.

Prior LLM-driven miners cast factor discovery as a regularized predictive-fit problem,
selecting the factor that maximizes a scalar effectiveness measure $\mathcal{L}$
(typically the Information Coefficient or a rank variant) against a novelty/complexity
penalty $\mathcal{R}$:
\begin{equation}
  f^\star \;=\; \arg\max_{f\in\Fspace}\;
  \mathcal{L}\!\big(f(\bm{X}),\,\bm{y}\big)\;-\;\lambda\,\mathcal{R}(f).
  \label{eq:frictionless}
\end{equation}
Equation~\eqref{eq:frictionless} is, with minor variation, the objective of the broader
IC-maximizing literature~\cite{han2026quantaalpha,shi2025alphajungle}. It is
\emph{frictionless}: the
effectiveness term $\mathcal{L}$ is a correlation between a signal and a forward return
and is silent on whether the signal can be realized as profit once trading is priced
in.

%----------------------------------------------------------------------
\subsection{The gap: frictionless prediction vs.\ tradeable performance}
\label{sub:gap}

A signal optimized purely for predictive correlation is not the same object as a
tradeable factor. Maximizing~\eqref{eq:frictionless} systematically rewards signals
whose predictive content is concentrated in small, illiquid names and whose exploitation
requires high portfolio turnover---precisely the signals whose apparent edge is consumed
by bid--ask spread, slippage, market impact, and the price drift incurred between the
moment a signal is computed and the moment an order actually fills. Two frictions are
routinely omitted from~\eqref{eq:frictionless} and are the focus of this work:
\begin{enumerate}[leftmargin=1.4em,itemsep=2pt]
  \item \textbf{Transaction cost.} Turnover is priced, not free. In the IC-maximizing
        literature it appears, if at all, as a soft post-hoc penalty rather than inside
        the reward.
  \item \textbf{Transaction latency.} A factor is computed on information available at
        decision time $t$, but no order fills at the decision-time price. The order
        fills after a nonzero latency, at a later---and adversarially different---price.
        The realized entry (buy) price is therefore a function of latency, not of the
        price the signal was computed on.
\end{enumerate}
We therefore replace the frictionless objective~\eqref{eq:frictionless} with a
\emph{net-of-cost, capacity-aware} objective in which both frictions are priced
\emph{inside} the quantity the miner optimizes, and evaluated by a deterministic engine
whose parameters the generation process cannot alter (\S\ref{sub:protocol}).

%----------------------------------------------------------------------
\subsection{Execution model: cost and latency}
\label{sub:exec}

\paragraph{Latency and the fill price.}
Let $\delta\ge 0$ be a fixed \emph{execution latency}, measured in trading periods,
between the decision time $t$ and the time an order is filled. A signal $\sign_t$
computed at the close of $t$ is executed at a fill price determined at $t+\delta$:
\begin{equation}
  P^{\mathrm{fill}}_{i,t} \;=\; \Phi\!\big(P_i;\,t,\delta\big),
  \label{eq:fill}
\end{equation}
where $\Phi$ is a fixed fill rule (e.g., the open of $t{+}1$ for $\delta{=}1$, or a
VWAP over $[t,\,t+\delta]$). The tradeable per-asset return realized between consecutive
fills,
\begin{equation}
  \tilde{y}_{i,t} \;=\; \frac{P^{\mathrm{fill}}_{i,t+1}}{P^{\mathrm{fill}}_{i,t}}-1,
  \label{eq:realized}
\end{equation}
differs from the frictionless label $\bm{y}_{t+1}$ whenever $\delta>0$: latency
determines the entry price and hence the return the strategy actually earns.

\paragraph{Portfolio construction.}
A fixed, deterministic map $g$ converts the composite
prediction~\eqref{eq:combiner} into portfolio weights
$\wt_t = g(\pred_t)\in\R^{N}$, \emph{not} an idealized equal-weight top decile. Weights
are signed---$w_{i,t}>0$ is a long and $w_{i,t}<0$ is a short position---so the
formulation considers both long and short selling; the canonical construction is a
rank-based, dollar-neutral long--short book that goes long the top of the cross-section
and short the bottom, with the gross-exposure (leverage) constraint
$\sum_i |w_{i,t}|\le 1$. The one-way \emph{turnover} at $t$ is
\begin{equation}
  \mathrm{TO}_t \;=\; \tfrac{1}{2}\textstyle\sum_{i=1}^{N}
  \big|\,w_{i,t}-w^{\mathrm{drift}}_{i,t}\,\big|,
  \label{eq:turnover}
\end{equation}
where $w^{\mathrm{drift}}_{i,t}$ is the pre-rebalance weight after market drift.

\paragraph{Cost and net return.}
Trading incurs three modeled cost components. Writing the per-asset traded fraction as
$\Delta w_{i,t}=w_{i,t}-w^{\mathrm{drift}}_{i,t}$, the per-period cost is
\begin{equation}
  c_t \;=\;
  \underbrace{\kappa_0\,\mathrm{TO}_t}_{\text{spread}+\text{fees}}
  \;+\;
  \underbrace{\kappa_1\!\sum_{i}\sigma_{i,t}\,\big|\Delta w_{i,t}\big|}_{\text{slippage}}
  \;+\;
  \underbrace{\kappa_2\!\sum_{i}\phi\!\big(\Delta w_{i,t};\,\mathrm{ADV}_{i,t}\big)}_{\text{market impact}}
  \;+\;
  \underbrace{\sum_{i}\beta_{i,t}\,\max(0,-w_{i,t})}_{\text{short borrow}},
  \label{eq:cost}
\end{equation}
with fixed coefficients $(\kappa_0,\kappa_1,\kappa_2)$ and a fixed, convex impact
function $\phi$ (e.g., a square-root law $\phi(\Delta w;\mathrm{ADV})\propto
|\Delta w|^{3/2}/\sqrt{\mathrm{ADV}}$). The first three terms are traded-size costs and
the fourth is a holding cost on the short leg: $\beta_{i,t}\ge 0$ is the per-period
borrow/financing rate for shorting asset $i$, charged on the short notional
$\max(0,-w_{i,t})$, which also encodes hard-to-borrow names (large or infinite
$\beta_{i,t}$) so that unshortable assets are penalized rather than assumed freely
sellable. The size-dependent terms are deliberately distinct: the
first is a size-independent charge proportional to turnover; the second is
\emph{slippage}, the temporary execution shortfall from crossing the book, taken linear
in traded size and scaled by the asset's return volatility $\sigma_{i,t}$; the third is
\emph{market impact}, which grows super-linearly in trade size relative to available
liquidity and is what makes the objective capacity-aware. The remaining, timing
component of slippage---price drift between the decision at $t$ and the fill at
$t+\delta$---is not modeled here but realized directly through the latency-aware fill
price of Eqs.~\eqref{eq:fill}--\eqref{eq:realized}, so it is not double-counted in $c_t$.
The \emph{net excess return} of the portfolio, computed on the latency-aware realized
returns~\eqref{eq:realized}, is
\begin{equation}
  r^{\mathrm{net}}_{t}
  \;=\; \wt_t^{\!\top}\tilde{\bm{y}}_{t}
  \;-\; r^{\mathrm{bench}}_{t}
  \;-\; c_t .
  \label{eq:netreturn}
\end{equation}
Equations~\eqref{eq:fill}--\eqref{eq:netreturn} constitute the execution model: the
single scalar that any downstream objective is built from is
$r^{\mathrm{net}}_t$, which prices turnover, slippage/impact, and latency
simultaneously.

%----------------------------------------------------------------------
\subsection{Multi-objective factor quality}
\label{sub:objectives}

We score a candidate factor $f$ on a vector of quality dimensions
$\bm{m}(f)=\big(m_1,\dots,m_M\big)$, extending the multi-dimensional evaluation of
formulaic miners~\cite{shi2025alphajungle} so that every economically meaningful
dimension is computed on the net-of-cost, latency-aware quantities of
\S\ref{sub:exec}. All statistics use $\bar{(\cdot)}$ for the temporal mean and
$\sigma(\cdot)$ for the temporal standard deviation over the evaluation window.

\begin{enumerate}[leftmargin=1.6em,itemsep=3pt]
\item \textbf{Effectiveness.} Predictive alignment via the Information Coefficient
      $\mathrm{IC}_t=\mathrm{corr}(\sign_t,\bm{y}_{t+1})$ and its rank analogue
      $\mathrm{RankIC}_t$, together with the realized economic edge, the
      \emph{net-of-cost} Information Ratio
      $\mathrm{IR}=\big(\overline{r^{\mathrm{net}}}/\sigma(r^{\mathrm{net}})\big)\sqrt{P}$,
      where $P$ is the number of periods per year.

\item \textbf{Annualized return.} The annualized net excess return
      $\mathrm{ARR}=\big(\prod_t (1+r^{\mathrm{net}}_t)\big)^{P/T_{\text{eval}}}-1$,
      reported as a headline economic objective alongside maximum drawdown (MDD) for
      risk context.

\item \textbf{Stability.} Temporal consistency of the signal, via the IC information
      ratios $\mathrm{ICIR}=\overline{\mathrm{IC}}/\sigma(\mathrm{IC})$ and
      $\mathrm{RankICIR}$, and the fraction of periods with positive IC. Stability
      penalizes signals whose edge is realized in a few episodes and is otherwise
      noise.

\item \textbf{Turnover.} Mean one-way turnover $\overline{\mathrm{TO}}$
      from~\eqref{eq:turnover}. Beyond feeding the cost term~\eqref{eq:cost}, turnover
      is a first-class capacity dimension: it bounds the size at which the factor can be
      deployed.

\item \textbf{Diversity.} Redundancy against the existing repository $\zoo$, measured
      as the maximum absolute cross-sectional correlation of the candidate's signal
      with any incumbent, $\rho_{\max}(f)=\max_{f'\in\zoo}\big|\mathrm{corr}(\sign^{f},\sign^{f'})\big|$.
      Low $\rho_{\max}$ favors orthogonal signals and directly combats factor crowding.

\item \textbf{Over-fitting risk.} A composite of (i) structural complexity
      $\mathrm{cx}(f)$---expression length / node count / free-parameter count of the
      factor's syntax tree, following the AST-based complexity control
      of~\cite{tang2025alphaagent}---and (ii) the in-sample-to-out-of-sample
      degradation $\Delta_{\mathrm{IS\to OOS}}(f)=m^{\mathrm{IS}}(f)-m^{\mathrm{OOS}}(f)$
      of the effectiveness metric. Higher complexity and larger degradation both raise
      the estimated risk that a backtested edge is spurious.

\item \textbf{Decay resistance (optional).} Persistence of predictive power across the
      out-of-sample horizon. Partition the OOS window into $K$ consecutive sub-periods
      and let $\mathrm{IC}^{(k)}$ be the mean IC in sub-period $k$. Fit the linear trend
      $\mathrm{IC}^{(k)}=a+b\,k+\varepsilon_k$; the \emph{decay slope} $b$ (and the
      persistence ratio $\mathrm{IC}^{(K)}/\mathrm{IC}^{(1)}$) summarize how quickly the
      edge erodes. We report a persistence-adjusted effectiveness
      $\mathrm{IC}^{\mathrm{pers}}=\overline{\mathrm{IC}}\cdot\big(1-\lambda_d\max(0,-\hat{b}_{\text{norm}})\big)$
      as a decay-aware summary. See Remark~\ref{rem:decay} on its status and scope.
\end{enumerate}

\paragraph{Adaptive scoring against the evolving repository.}
Admissibility and scoring are distinct steps. The gates of~\eqref{eq:feasible} decide
\emph{whether} a factor enters the feasible set $\feasible$, using the fixed, absolute
floors above; once a factor is admissible, we must still decide \emph{how good} it is
relative to the others, and here fixed cutoffs are the wrong tool. Because the
effective-alpha repository $\zoo$ grows as mining proceeds, a scoring bar that is
informative early becomes trivial late; we therefore score each admissible factor's
dimensions by their \emph{relative rank} within the current repository,
following~\cite{shi2025alphajungle}.
For a dimension metric $m$ (oriented so that larger is better), the rank of a candidate
against the repository is
\begin{equation}
  R(f, m, \zoo) \;=\; \frac{1}{|\zoo|}\sum_{f'\in\zoo}\mathbb{1}\!\big[\,m(f) < m(f')\,\big],
  \label{eq:relrank}
\end{equation}
where $\mathbb{1}[\cdot]$ is the indicator function, and the corresponding dimension
score is the percentile complement $e_i(f)=1-R(f,m_i,\zoo)\in[0,1]$, so a factor that
outperforms more incumbents scores higher. This gives an adaptive criterion whose bar
rises automatically as $\zoo$ improves, avoiding fixed thresholds that are too stringent
at the start or too lenient once many strong factors have accumulated. Over-fitting risk
(item~6) is the exception: it is assessed from the factor's structure and refinement
history rather than by repository rank. The scalar utility $\mathcal{U}$ of
\S\ref{sub:statement} aggregates these per-dimension scores $e_i(f)$. Note the two roles
the repository plays are distinct: for \emph{combination} (Eq.~\eqref{eq:combiner}) the
library is frozen to fit $\combiner$, while for \emph{scoring} (Eq.~\eqref{eq:relrank})
it is the moving reference set against which each new candidate is ranked; the same
accumulated set of effective alphas underlies both.

\begin{remark}[On ``decay-resistant IC'']
\label{rem:decay}
Decay resistance is a well-motivated \emph{property}~\cite{tang2025alphaagent}: factors
lose predictive power as regimes shift and crowding sets in, and a net-of-cost objective
is by construction non-stationary. It is not, however, a standard closed-form metric;
in prior work it is achieved through regularization and demonstrated \emph{empirically}
as IC/IR persistence over time. We therefore treat it as an out-of-sample persistence
diagnostic (the decay slope $b$ and $\mathrm{IC}^{\mathrm{pers}}$ above) rather than a
new headline statistic. It is partially---but not fully---subsumed by Stability
(which measures IC \emph{variance}) and Over-fitting risk (which measures the IS$\to$OOS
\emph{level gap}): neither captures a monotone \emph{temporal trend} within the OOS
window, which is the distinct signal $b$ provides. We include it as an optional
diagnostic and defer a formal significance treatment of alpha decay to the falsification
component discussed in \S\ref{sub:statement}.
\end{remark}

%----------------------------------------------------------------------
\subsection{Strict deterministic backtesting protocol}
\label{sub:protocol}

The objective above is only as trustworthy as the procedure that computes it. We
enforce a strict separation of concerns, in the spirit of the deterministic evaluation
engine of~\cite{huang2026hypotheses}: the generation process proposes factors, while a
fixed evaluation engine---and only the engine---assigns scores.

Let the \emph{protocol}
$\proto=\big(\mathcal{D}_{\mathrm{tr}},\mathcal{D}_{\mathrm{va}},\mathcal{D}_{\mathrm{te}};\,
\gates;\,\kappa_0,\kappa_1,\kappa_2,\phi,\beta;\,\delta,\Phi;\,g\big)$
collect every parameter of evaluation: the frozen train/validation/test splits, the
selection gates $\gates$, the cost model, the latency and fill rule, and the portfolio
map. Given $\proto$, a \emph{deterministic evaluation operator} $\evalop_\proto$ maps any
factor to its metric vector,
\begin{equation}
  \bm{m}(f) \;=\; \evalop_\proto(f),
  \label{eq:evalop}
\end{equation}
computed strictly on point-in-time data with no look-ahead. The engine guarantees:

\begin{property}[Immutability]
\label{prop:immut}
$\proto$ is fixed before search and is a constant of the environment. The
generation process can propose any $f\in\Fspace$ but cannot alter any element of
$\proto$---it cannot move a window boundary, relax a gate threshold, re-price a cost, or
change the latency. Reported metrics thus cannot be inflated by tuning the protocol to
the factor---the reuse of evaluation windows during search being a documented source of
performance overestimation~\cite{deng2025autoquant}.
\end{property}

\begin{property}[Determinism and reproducibility]
\label{prop:det}
$\evalop_\proto$ is a pure function of $(f,\proto)$: identical inputs yield identical
$\bm{m}(f)$, so every trial is exactly reproducible and auditable.
\end{property}

\begin{property}[Point-in-time integrity]
\label{prop:pit}
All features, labels~\eqref{eq:realized}, fills~\eqref{eq:fill}, and
costs~\eqref{eq:cost} use only information available at or before their timestamp;
the test split $\mathcal{D}_{\mathrm{te}}$ is never touched during selection.
\end{property}

The \emph{selection gates} $\gates$ act as hard feasibility constraints, evaluated on
the training/validation splits only:
\begin{equation}
  \feasible \;=\; \Big\{\, f\in\Fspace \;:\;
  \mathrm{RankIC}(f)\ge \gamma_{\mathrm{ic}},\;
  \mathrm{RankICIR}(f)\ge \gamma_{\mathrm{ir}},\;
  \overline{\mathrm{TO}}(f)\le \tau_{\max},\;
  \rho_{\max}(f)\le \bar{\rho},\;
  \mathrm{cx}(f)\le \gamma_{\mathrm{cx}}
  \,\Big\},
  \label{eq:feasible}
\end{equation}
with all thresholds $(\gamma_{\mathrm{ic}},\gamma_{\mathrm{ir}},\tau_{\max},\bar{\rho},
\gamma_{\mathrm{cx}})\in\proto$ frozen. Gates encode the tradeability and
diversity requirements as \emph{admissibility}, not as soft preferences: they answer the
binary question of whether a factor can be traded at all, and their thresholds are
absolute physical floors---a turnover cap or a redundancy ceiling is a property of the
market and the existing book, not a bar that should drift with search progress. This is a
separate question from how the admissible factors are then \emph{scored} and ranked
against one another, which we turn to next.

%----------------------------------------------------------------------
\subsection{Problem statement}
\label{sub:statement}

We can now state the problem. Given a frozen protocol $\proto$ and its deterministic
evaluation operator $\evalop_\proto$, we seek factors that maximize a net-of-cost
economic utility over the feasible set defined by the frozen gates:
\begin{equation}
  \boxed{\;
  f^\star \;=\; \arg\max_{f\in\feasible}\;
  \mathcal{U}\!\big(\bm{m}(f)\big),
  \qquad \bm{m}(f)=\evalop_\proto(f),
  \;}
  \label{eq:problem}
\end{equation}
where $\mathcal{U}$ is a net-of-cost economic objective---at minimum the net Information
Ratio / ARR of~\eqref{eq:netreturn}, optionally a fixed scalarization
$\mathcal{U}(\bm{m})=\sum_j \omega_j \hat{m}_j$ with $\hat{m}_j \equiv e_j(f)$, i.e.\ the
normalized quality dimensions of \S\ref{sub:objectives} are exactly the
repository-relative dimension scores $e_j(f)=1-R(f,m_j,\zoo)$ of
Eq.~\eqref{eq:relrank}, combined with fixed weights $\bm{\omega}\in\proto$ to rank
candidates and give dimension-targeted feedback to the generation process. Because
the repository $\zoo$ is built incrementally, both the diversity term $\rho_{\max}$
in~\eqref{eq:feasible} and the relative ranks in~\eqref{eq:relrank} are evaluated against
the current $\zoo$, so~\eqref{eq:problem} defines the marginal contribution of each factor
to a jointly non-redundant, capacity-robust portfolio.

Two properties distinguish~\eqref{eq:problem} from the frictionless
objective~\eqref{eq:frictionless}. First, \emph{costs and latency are inside the
objective}: $\bm{m}(f)$ is computed from $r^{\mathrm{net}}$, so the miner is rewarded for
capacity-robust structure rather than for illiquid, high-turnover mirages. Second, the
objective is \emph{honest by construction}: by Properties~\ref{prop:immut}--\ref{prop:pit},
the score of every candidate is a deterministic, reproducible, look-ahead-free function
of a protocol the generator cannot bend to its own factors.

\paragraph{Scope and roadmap.}
This work is the first component of a larger system and is deliberately scoped to the
objective and its enforcement: the factor-generation process is held fixed and we change
only what it optimizes and how that is measured. Concretely, the evaluation protocol
$\proto$---markets, data splits, feature set, and the combination model
$\combiner$---is instantiated with the configuration of a published IC-maximizing
baseline~\cite{han2026quantaalpha} and held fixed, so that the baseline's reported
results are directly comparable and the effect of the objective change can be isolated
without re-running the baseline pipeline. Subsequent components---a falsification critic that supplies a
selection-bias-corrected significance test over the complete population of trials
(e.g., a deflated Sharpe ratio~\cite{bailey2014dsr} under multiple-testing
correction~\cite{harvey2016cross}), and a
continually refreshed, decay-monitored factor library with per-factor provenance for
deployment---build on the reproducible trial record that
Properties~\ref{prop:immut}--\ref{prop:det} make available, and are outside the scope of
this section.

%======================================================================
\begin{thebibliography}{9}
\small
\bibitem{han2026quantaalpha}
J.~Han, S.~Zhang, W.~Li, et al.
\newblock QuantaAlpha: An Evolutionary Framework for LLM-Driven Alpha Mining.
\newblock \emph{arXiv:2602.07085}, 2026.

\bibitem{shi2025alphajungle}
Y.~Shi, Y.~Duan, and J.~Li.
\newblock Navigating the Alpha Jungle: An LLM-Powered MCTS Framework for Formulaic
Factor Mining.
\newblock \emph{arXiv:2505.11122}, 2025.

\bibitem{tang2025alphaagent}
Z.~Tang, et al.
\newblock AlphaAgent: LLM-Driven Alpha Mining with Regularized Exploration to
Counteract Alpha Decay.
\newblock \emph{Proc.\ ACM SIGKDD (KDD~'25)}; arXiv:2502.16789, 2025.

\bibitem{huang2026hypotheses}
Y.~Huang, Z.~Fan, K.~Hu, and Y.~Ye.
\newblock From Hypotheses to Factors: Constrained LLM Agents in Cryptocurrency Markets.
\newblock \emph{arXiv:2604.26747}, 2026.

\bibitem{deng2025autoquant}
K.~Deng.
\newblock AutoQuant: An Auditable Expert-System Framework for Execution-Constrained
Auto-Tuning in Cryptocurrency Perpetual Futures.
\newblock \emph{arXiv:2512.22476}, 2025.

\bibitem{bailey2014dsr}
D.~H. Bailey and M.~L\'opez de Prado.
\newblock The Deflated Sharpe Ratio: Correcting for Selection Bias, Backtest
Overfitting, and Non-Normality.
\newblock \emph{J.\ Portfolio Management}, 40(5), 2014.

\bibitem{harvey2016cross}
C.~R. Harvey, Y.~Liu, and H.~Zhu.
\newblock $\ldots$ and the Cross-Section of Expected Returns.
\newblock \emph{Review of Financial Studies}, 29(1), 2016.
\end{thebibliography}

\end{document}
